\subsection{Implementazione della ricerca Negamax}

Per affiancare alla selezione diretta delle mosse una componente di ricerca, è stato implementato un agente basato sull'algoritmo Negamax con potatura alpha-beta. Questa modalità non sostituisce la rete TRM, ma la utilizza come supporto per individuare e ordinare le mosse da esplorare.

La TRM continua infatti a generare una distribuzione di probabilità sulle azioni possibili, mentre Negamax utilizza tali informazioni per costruire un albero di ricerca composto da sequenze alternative di mosse. L'obiettivo è valutare non soltanto la mossa immediatamente più probabile, ma anche le possibili risposte dell'avversario e le conseguenze delle proprie scelte.

\subsubsection{Generazione e ordinamento delle mosse candidate}

A ogni nodo dell'albero, l'agente interroga la policy TRM per ottenere le probabilità associate alle azioni disponibili. Le mosse illegali vengono escluse utilizzando la maschera fornita dall'ambiente di gioco.

Le mosse legali vengono quindi ordinate in base alla probabilità assegnata dalla rete e vengono selezionate le prime $K$, dove $K$ rappresenta il numero massimo di candidate considerate per ciascun nodo. Questo limite consente di contenere il numero di rami esplorati, evitando di dover analizzare tutte le mosse possibili in ogni posizione.

La selezione non si basa esclusivamente sulle indicazioni della rete. Vengono infatti aggiunte alcune mosse tatticamente rilevanti, anche quando la loro probabilità è inferiore rispetto a quella delle candidate selezionate inizialmente. In particolare, vengono considerate le libertà delle catene di pietre che si trovano in atari, così da includere le possibili mosse di cattura o di salvataggio.

Anche la mossa PASS può essere inclusa tra le candidate: viene considerata quando rientra tra le azioni con probabilità più elevata oppure quando l'avversario ha già effettuato un PASS, situazione nella quale è necessario valutare anche la possibilità di terminare la partita.

Le mosse candidate vengono infine ordinate, dando precedenza a quelle tattiche e utilizzando successivamente la probabilità della policy come criterio di ordinamento. Questo approccio permette di combinare le indicazioni apprese dalla rete con alcune informazioni tattiche ricavate direttamente dalla configurazione della scacchiera.

\subsubsection{Ricerca a profondità limitata e potatura alpha-beta}

Una volta individuate le mosse candidate, l'agente esplora ricorsivamente le possibili sequenze di gioco attraverso l'algoritmo Negamax. La ricerca procede alternando le mosse dei due giocatori e valutando ogni posizione dal punto di vista del giocatore che deve effettuare la mossa.

La profondità massima della ricerca è configurabile e indica il numero di mosse consecutive considerate nell'albero. Ad esempio, una profondità pari a due permette di valutare una possibile mossa del giocatore corrente e la successiva risposta dell'avversario.

L'algoritmo Negamax sfrutta il fatto che, in un gioco a somma zero, il valore di una posizione per un giocatore corrisponde all'opposto del valore della stessa posizione per l'avversario. Di conseguenza, a ogni livello ricorsivo il valore restituito dalla chiamata successiva viene invertito di segno.

Per ridurre il numero di posizioni da esplorare viene utilizzata la potatura alpha-beta. Durante la ricerca vengono mantenuti due limiti, $\alpha$ e $\beta$, che rappresentano rispettivamente il miglior valore già garantito al giocatore corrente e il limite oltre il quale il giocatore avversario non permetterebbe di proseguire lungo quel ramo. Quando $\alpha \geq \beta$, le ulteriori mosse del ramo possono essere escluse dalla ricerca, poiché non potrebbero influenzare la decisione finale.

Prima di esplorare ciascuna mossa, viene inoltre creata una copia indipendente dell'ambiente di gioco. In questo modo ogni ramo dell'albero può modificare la propria configurazione della scacchiera senza alterare lo stato della partita originale o quello degli altri rami.

\subsubsection{Funzione euristica di valutazione}

Quando la ricerca raggiunge la profondità massima senza incontrare una condizione terminale, la posizione viene valutata attraverso una funzione euristica. Questa valutazione non è prodotta da una rete neurale dedicata, ma viene calcolata direttamente a partire dallo stato della partita.

La funzione considera principalmente due aspetti:

\begin{itemize}
    \item \textbf{Differenza del punteggio territoriale}: viene calcolata la differenza tra il punteggio del Nero e quello del Bianco, secondo il sistema di conteggio dell'area utilizzato dall'ambiente di gioco.
    \item \textbf{Sicurezza dei gruppi}: vengono analizzate le libertà delle catene di pietre presenti sulla scacchiera, applicando una penalità ai gruppi più vulnerabili. In particolare, la penalità è maggiore per i gruppi in atari o con un numero ridotto di libertà.
\end{itemize}

Il contributo dei due aspetti viene combinato in un unico valore, successivamente adattato al punto di vista del giocatore corrente. In questo modo, una posizione viene valutata positivamente quando presenta un vantaggio territoriale e una configurazione tatticamente favorevole per il giocatore considerato.

È importante sottolineare che questa funzione rappresenta una stima euristica della posizione e non una valutazione appresa durante l'addestramento. La TRM, infatti, è stata progettata e addestrata per generare una policy, ossia una distribuzione di probabilità sulle mosse, e non per stimare direttamente il valore di una posizione o la probabilità di vittoria.

\subsubsection{Gestione delle condizioni terminali}

Durante la ricerca viene verificato a ogni nodo se la partita è terminata. Nel caso in cui venga raggiunta una posizione terminale, la valutazione viene effettuata utilizzando il punteggio effettivo della partita, anziché la funzione euristica.

Il valore terminale viene calcolato a partire dalla differenza tra il punteggio del Nero e quello del Bianco, adattandolo al giocatore di riferimento. Per distinguere chiaramente un risultato definitivo da una semplice stima euristica, viene inoltre aggiunto un bonus di ampiezza elevata al valore delle vittorie e delle sconfitte.

In questo modo, quando la ricerca individua una sequenza che conduce a una vittoria o a una sconfitta, il relativo risultato assume priorità rispetto alle valutazioni delle posizioni non terminali.

Al termine della ricerca, l'agente restituisce la prima mossa della sequenza individuata come migliore. L'ambiente di gioco applica quindi tale mossa e aggiorna lo stato effettivo della partita.

\paragraph{Considerazioni sull'integrazione con la TRM}

L'introduzione di Negamax costituisce un'estensione del processo di inferenza e non comporta modifiche alla rete neurale. I pesi della TRM rimangono invariati e vengono utilizzati per generare le probabilità delle mosse candidate durante la ricerca.

La componente Negamax opera quindi al di sopra della policy esistente, combinando le indicazioni del modello con l'esplorazione di sequenze di gioco, la valutazione euristica delle posizioni e la potatura alpha-beta. Non è pertanto necessario eseguire un nuovo addestramento per utilizzare questa modalità: è sufficiente caricare il checkpoint della TRM già addestrata e configurare i parametri della ricerca, come la profondità massima e il numero di mosse candidate.